

\medskip
\noindent\SGAInlineRefTarget{sga3:iii:paragraph:4-5-0}\textbf{4.5.0.}
Retain the notation of \SGARefLink{sga3:iii:corollary:4-3-1}{4.3.1} and \SGARefLink{sga3:iii:corollary:4-4}{4.4}, and make the following
transformations. The quasi-coherent \(\mathscr O_{X_{\mathscr J}}\)-module
\[
 \mathscr I_{Y_{\mathscr J}}/\mathscr I_{Y_{\mathscr J}}^2
\]
is annihilated by \(\mathscr I_{Y_{\mathscr J}}\), and is therefore the
direct image of a quasi-coherent \(\mathscr O_{Y_{\mathscr J}}\)-module,
denoted by \(\mathscr N_{Y_{\mathscr J}/X_{\mathscr J}}\), and called the
\emph{conormal sheaf} of \(Y_{\mathscr J}\) in
\(X_{\mathscr J}\).
Since
\(\mathscr J\otimes_{\mathscr O_{S_{\mathscr J}}}
\mathscr O_{Y_{\mathscr J}}\) is annihilated by
\(\mathscr I_{Y_{\mathscr J}}\), the sheaf of commutative groups
\(\mathscr A_{\mathscr J}\) of \SGARefLink{sga3:iii:corollary:4-3-1}{4.3.1} identifies with
\[
 \begin{split}
 &\operatorname{Hom}_{\mathscr O_{Y_{\mathscr J}}}
 \left(
  \mathscr I_{Y_{\mathscr J}}/\mathscr I_{Y_{\mathscr J}}^2,
  \mathscr J\otimes_{\mathscr O_{S_{\mathscr J}}}
   \mathscr O_{Y_{\mathscr J}}
 \right) \\
 &\qquad={}
 \operatorname{Hom}_{\mathscr O_{Y_{\mathscr J}}}
 \left(
  \mathscr N_{Y_{\mathscr J}/X_{\mathscr J}},
  \mathscr J\otimes_{\mathscr O_{S_{\mathscr J}}}
   \mathscr O_{Y_{\mathscr J}}
 \right),
 \end{split}
\]
and hence, for every \(i\geq0\),
\[
 H^i(X_{\mathscr J},\mathscr A_{\mathscr J})
 =H^i\!\left(
  Y_{\mathscr J},
  \operatorname{Hom}_{\mathscr O_{Y_{\mathscr J}}}
  \left(
   \mathscr N_{Y_{\mathscr J}/X_{\mathscr J}},
   \mathscr J\otimes_{\mathscr O_{S_{\mathscr J}}}
    \mathscr O_{Y_{\mathscr J}}
  \right)
 \right).
\]

\footnote{Editor's note: The discussion that follows has been expanded.}
One may now drop the hypothesis that \(Y\) be closed, as follows. First
observe that every open subset \(U_{\mathscr J}\) of
\(X_{\mathscr J}\) is obtained by base change from the open subscheme
\(U\) of \(X\) having the same underlying topological space as
\(U_{\mathscr J}\). Let \(Y_{\mathscr J}\) now be a closed subscheme of
\(U_{\mathscr J}\), flat over \(S_{\mathscr J}\), and let
\(\mathscr I_{Y_{\mathscr J}}\) be the quasi-coherent ideal of
\(\mathscr O_{U_{\mathscr J}}\) defining \(Y_{\mathscr J}\). If
\(Y_{\mathscr J}\) lifts to a subscheme \(Y\) of \(X\), then \(Y\),
having the same underlying topological space as \(Y_{\mathscr J}\), is a
closed subscheme of \(U\). Consequently, the obstruction to lifting
\(Y_{\mathscr J}\) to a subscheme, flat over \(S\), of \(X\) or of
\(U\) is the same; it lies in
\[
 H^1\!\left(
  Y_{\mathscr J},
  \operatorname{Hom}_{\mathscr O_{Y_{\mathscr J}}}
  \left(
   \mathscr N_{Y_{\mathscr J}/X_{\mathscr J}},
   \mathscr J\otimes_{\mathscr O_{S_{\mathscr J}}}
    \mathscr O_{Y_{\mathscr J}}
  \right)
 \right).
\]

Finally, return to the notation of \S0. Let \(\mathscr I\) be a
quasi-coherent ideal of \(\mathscr O_S\) such that
\(\mathscr J\subset\mathscr I\) and \(\mathscr I\mathscr J=0\), and let
\(S_0\) be the closed subscheme of \(S_{\mathscr J}\) defined by
\(\mathscr I\). For every \(S\)-scheme \(Z\), write
\[
 Z_{\mathscr J}=Z\times_S S_{\mathscr J},
 \qquad Z_0=Z\times_S S_0.
\]
Since \(\mathscr J\) is annihilated by \(\mathscr I\), the notation of
\SGARefLink{sga3:iii:corollary:4-4}{4.4} gives
\[
 \mathscr J\otimes_{\mathscr O_{S_{\mathscr J}}}
  \mathscr O_{Y_{\mathscr J}}
 =\mathscr J\otimes_{\mathscr O_{S_0}}\mathscr O_{Y_0},
\]
and
\[
 \begin{split}
 &\operatorname{Hom}_{\mathscr O_{Y_{\mathscr J}}}
 \left(
  \mathscr N_{Y_{\mathscr J}/X_{\mathscr J}},
  \mathscr J\otimes_{\mathscr O_{S_{\mathscr J}}}
   \mathscr O_{Y_{\mathscr J}}
 \right) \\
 &\qquad={}
 \operatorname{Hom}_{\mathscr O_{Y_0}}
 \left(
  \mathscr N_{Y_{\mathscr J}/X_{\mathscr J}}
   \otimes_{\mathscr O_{Y_{\mathscr J}}}\mathscr O_{Y_0},
  \mathscr J\otimes_{\mathscr O_{S_0}}\mathscr O_{Y_0}
 \right),
 \end{split}
\]
and similarly for the other terms. Thus one obtains:

\medskip
\noindent\SGAInlineRefTarget{sga3:iii:proposition:4-5}\textbf{Proposition 4.5.}
\textit{Let \(S\) be a scheme, and let \(S_0\) and
\(S_{\mathscr J}\) be the closed subschemes defined by quasi-coherent
ideals \(\mathscr I\) and \(\mathscr J\), where
\(\mathscr I\supset\mathscr J\) and
\(\mathscr I\mathscr J=0\). Let \(X\) be an \(S\)-scheme and let
\(Y_{\mathscr J}\) be a subscheme of \(X_{\mathscr J}\), flat over
\(S_{\mathscr J}\). Define the \(\mathscr O_{Y_0}\)-module
\[
 \mathscr A_0=
 \operatorname{Hom}_{\mathscr O_{Y_0}}
 \left(
  \mathscr N_{Y_{\mathscr J}/X_{\mathscr J}}
   \otimes_{\mathscr O_{Y_{\mathscr J}}}\mathscr O_{Y_0},
  \mathscr J\otimes_{\mathscr O_{S_0}}\mathscr O_{Y_0}
 \right).
\]
Then:}
\begin{enumerate}
\item[\SGAInlineRefTarget{sga3:iii:proposition:4-5:item:i}(i)] \textit{A subscheme \(Y\) of \(X\), reducing to
  \(Y_{\mathscr J}\) and flat over \(S\), exists if and only if the
  following two conditions hold:}
  \begin{enumerate}
  \item[\SGAInlineRefTarget{sga3:iii:proposition:4-5:item:a}(a)] \textit{Such a \(Y\) exists locally on \(X\).}
  \item[\SGAInlineRefTarget{sga3:iii:proposition:4-5:item:b}(b)] \textit{A certain obstruction in
    \(R^1\Gamma(Y_0,\mathscr A_0)\) vanishes.}
  \end{enumerate}
\item[\SGAInlineRefTarget{sga3:iii:proposition:4-5:item:ii}(ii)] \textit{Under these conditions, the set of subschemes \(Y\)
  satisfying the stated requirements is principal homogeneous under the
  abelian group \(\Gamma(Y_0,\mathscr A_0)\).}
\end{enumerate}

\medskip
\noindent\SGAInlineRefTarget{sga3:iii:remark:4-5-1}\textbf{Remark 4.5.1.}\footnote{Editor's note: This remark has
been placed here; it replaces \SGARefLink{sga3:iii:paragraph:4-7}{Remark 4.7} of the original.}
It follows from \SGARefLink{sga3:iii:proposition:4-5:item:ii}{4.5(ii)} that for every pair \((Y,Y')\) of
subschemes of \(X\), flat over \(S\) and reducing to
\(Y_{\mathscr J}\), there is a ``deviation''
\[
 \begin{split}
 d(Y',Y)&\in\Gamma(Y_0,\mathscr A_0)\\
 &=\operatorname{Hom}_{\mathscr O_{Y_0}}
 \left(
  \mathscr N_{Y_{\mathscr J}/X_{\mathscr J}}
   \otimes_{\mathscr O_{Y_{\mathscr J}}}\mathscr O_{Y_0},
  \mathscr J\otimes_{\mathscr O_{S_0}}\mathscr O_{Y_0}
 \right).
 \end{split}
\]
With the sign convention adopted in \SGARefLink{sga3:iii:complement:4-4-1}{4.4.1}, \(d(Y',Y)\) corresponds to
the morphism of \(\mathscr O_X\)-modules
\[
 \mathscr I_{Y'}\lhook\joinrel\longrightarrow
 \mathscr O_X\longrightarrow\mathscr O_Y,
\]
which takes values in
\(\mathscr J\mathscr O_Y\simeq
\mathscr J\otimes_{\mathscr O_{S_0}}\mathscr O_{Y_0}\), and factors
through
\(\mathscr I_{Y'_{\mathscr J}}=\mathscr I_{Y_{\mathscr J}}\), then
through \(\mathscr N_{Y_{\mathscr J}/X_{\mathscr J}}\).

\medskip
\noindent\SGAInlineRefTarget{sga3:iii:remark:4-6}\textbf{Remark 4.6.}\footnote{Editor's note: We have retained,
for reference, \SGARefLink{sga3:iii:remark:4-6}{Remark 4.6} of the original, which did not include the
definition of ``locally complete intersection.'' It is followed by the
``correct'' definition, taken from SGA 6, VII 1.4 (and replacing that of
EGA IV\textsubscript{4}, 16.9.2), and by proofs of the three assertions
made in the remark.}
\textit{If \(X\) is flat over \(S\) and \(Y_{\mathscr J}\) is locally
complete intersection in \(X_{\mathscr J}\), then condition \SGARefLink{sga3:iii:proposition:4-5:item:a}{(a)} is
always satisfied, and every \(Y\) flat over \(S\) lifting
\(Y_{\mathscr J}\) is locally complete intersection in \(X\). If,
moreover, \(Y_0\) is affine, condition \SGARefLink{sga3:iii:proposition:4-5:item:b}{(b)} is also satisfied.}

\medskip
\noindent\SGAInlineRefTarget{sga3:iii:definition:4-6-1}\textbf{Definition 4.6.1.}
\textit{(cf. SGA 6, VII 1.1) Let \(B\) be a commutative ring, let
\(f:E\to B\) be a \(B\)-linear morphism, where \(E\) is a free
\(B\)-module of finite rank \(d\), and let \(I=f(E)\). If a basis of
\(E\) is chosen, \(f\) is given by a \(d\)-tuple
\((f_1,\ldots,f_d)\) of elements of \(B\), and \(I\) is the ideal
generated by the \(f_i\). The \emph{Koszul complex} \(K_\bullet(f)\)
is the graded \(B\)-module \(\bigwedge_B^\bullet E\), equipped with the
differential of degree \(-1\)}
\[
 x_1\wedge\cdots\wedge x_i
 \longmapsto
 \sum_{j=1}^i(-1)^{j-1}f(x_j)
 x_1\wedge\cdots\wedge\widehat{x_j}\wedge\cdots\wedge x_i.
\]
Thus there is an augmented chain complex, with \(B/I\) in degree
\(-1\),
\[
 \cdots\longrightarrow\bigwedge^2 E\longrightarrow E
 \xrightarrow{f}B\longrightarrow B/I\longrightarrow0,
\]
which is exact in degree \(0\) by definition, since \(I=f(E)\). The
morphism \(f\) is said to be \emph{regular} if \(K_\bullet(f)\) is
acyclic in degrees \(>0\), that is, if the augmented complex above is a
resolution of \(C=B/I\).

In that case, the proof of SGA 6, VII 1.2(b) shows that the
\(C\)-modules \(I^n/I^{n+1}\), for \(n\in\mathbb N\), are free, with
\(I/I^2\) of rank \(d\).

\medskip
\noindent\SGAInlineRefTarget{sga3:iii:definition:4-6-2}\textbf{Definition 4.6.2.}
\textit{(cf. SGA 6, VII 1.4) Let \(X\) be a scheme, \(Y\) a subscheme,
and \(U\) an open subset of \(X\) such that \(Y\) is a closed subscheme
of \(U\), defined by the quasi-coherent ideal \(\mathscr I_Y\).}

The subscheme \(Y\) is said to be \emph{locally complete intersection}
in \(X\) if \(Y\hookrightarrow X\) is a regular immersion in the sense
of SGA 6, VII 1.4; that is, if for every \(y\in Y\) there exist an affine
open neighborhood \(V\) of \(y\) in \(U\), a finite free
\(\mathscr O_V\)-module \(\mathscr E\), and a regular morphism
\(f:\mathscr E\to\mathscr O_V\) with image
\(\mathscr I_Y|_V\), such that \(K_\bullet(f)\) is a resolution of
\(\mathscr O_{Y\cap V}\).

This implies that the immersion \(Y\hookrightarrow X\) is locally of
finite presentation and, by \SGARefLink{sga3:iii:definition:4-6-1}{4.6.1}, that the conormal sheaf
\[
 \mathscr N_{Y/X}=\mathscr I_Y/\mathscr I_Y^2
\]
is a finite locally free \(\mathscr O_Y\)-module.

\medskip
\noindent\SGAInlineRefTarget{sga3:iii:lemma:4-6-3}\textbf{Lemma 4.6.3.}\footnote{\label{fnref:III-88}Editor's note: To prove the
results stated in \SGARefLink{sga3:iii:remark:4-6}{Remark 4.6}, \SGARefLink{sga3:iii:lemma:4-6-3}{Lemmas 4.6.3} and \SGARefLink{sga3:iii:lemma:4-6-4}{4.6.4}, \SGARefLink{sga3:iii:proposition:4-6-5}{Proposition 4.6.5},
and \SGARefLink{sga3:iii:remark:4-6-6}{Remark 4.6.6} have been added.}
\textit{Let \(A\) be a ring, let \(J\) be a square-zero ideal of
\(A\), put \(\overline A=A/J\), let \(B\) be a flat \(A\)-algebra, let
\(E\) be a free \(B\)-module of finite rank, and let \(f:E\to B\) be a
morphism of \(B\)-modules. Suppose that the morphism
\[
 g:\overline E=E\otimes_A\overline A
 \longrightarrow\overline B=B\otimes_A\overline A
\]
induced by \(f\) is regular and that
\(\overline B/g(\overline E)\) is flat over \(\overline A\). Then \(f\)
is regular and \(B/f(E)\) is flat over \(A\).}

\noindent\textit{Proof.}
Put
\[
 C=B/f(E),
 \qquad
 \overline C=C\otimes_A\overline A
 =\overline B/g(\overline E).
\]
First, the \(\bigwedge_B^i E\) are free \(B\)-modules, and hence flat
\(A\)-modules, since \(B\) is flat over \(A\). Since
\[
 \bigwedge_B^\bullet E\otimes_A\overline A
 \simeq\bigwedge_{\overline B}^\bullet\overline E,
\]
there is an exact sequence of complexes
\[
 0\longrightarrow
 J\otimes_A\bigwedge_B^\bullet E
 \longrightarrow\bigwedge_B^\bullet E
 \longrightarrow\bigwedge_{\overline B}^\bullet\overline E
 \longrightarrow0.
\]
Moreover, since \(J^2=0\), for every \(A\)-module \(M\) one has
\[
 J\otimes_A M
 \simeq J\otimes_{\overline A}(\overline A\otimes_A M).
\]
Writing the augmentation arrows as dashed arrows, and denoting the rank
of \(E\) by \(d\), one obtains the following bicomplex, whose rows are
exact.

\begin{center}
 \SGAThreeNativeInput[0.78\textwidth]
   {native_diagrams/exp3/diagram_038_exp3_p150_koszul_bicomplex.tex}
\end{center}

The right and left columns are also exact, since \(K_\bullet(g)\) is a
resolution of \(\overline C\), and \(\overline C\) is flat over
\(\overline A\). The long exact homology sequence associated with the
exact sequence of unaugmented complexes therefore shows that
\(K_\bullet(f)\) is acyclic in degrees \(>0\), and in degree \(0\) gives
an exact sequence
\[
 0\longrightarrow J\otimes_{\overline A}\overline C
 \longrightarrow C\longrightarrow\overline C\longrightarrow0.
\]
It follows from the ``fundamental criterion of flatness'' that \(C\) is
flat over \(A\) (cf. [BAC], \S III.5, Theorem 1).
\hfill\(\square\)

\medskip
\noindent\SGAInlineRefTarget{sga3:iii:lemma:4-6-4}\textbf{Lemma 4.6.4.}\textsuperscript{\normalfont\ref{fnref:III-88}}
\textit{Let \(A\) be a commutative ring, let \(J\) be a nilpotent ideal,
and let \(N\subset M\) be \(A\)-modules such that \(M/N\) is flat over
\(A\). If \(x_1,\ldots,x_n\) are elements of \(N\) whose images generate
the image \(\overline N\) of \(N\) in \(M/JM\), then they generate
\(N\).}

Indeed, let \(N'\) be the submodule of \(N\) generated by the \(x_i\),
and put \(Q=N/N'\). The morphism
\(N'\otimes_A(A/J)\to\overline N\) is surjective. On the other hand,
since \(M/N\) is flat over \(A\), the morphism
\(N\otimes_A(A/J)\to\overline N\) is bijective. Thus
\(Q\otimes_A(A/J)=0\), whence \(Q=0\) by the ``nilpotent Nakayama
lemma'':
\[
 Q=JQ=J^2Q=\cdots=0.
\]

We can now prove the following proposition.

\medskip
\noindent\SGAInlineRefTarget{sga3:iii:proposition:4-6-5}\textbf{Proposition 4.6.5.}\textsuperscript{\normalfont\ref{fnref:III-88}}
\textit{Let \(S,\mathscr I,\mathscr J\), and
\(X,Y_{\mathscr J}\) be as in \SGARefLink{sga3:iii:proposition:4-5}{4.5}. Suppose in addition that \(X\) is
flat over \(S\) and that \(Y_{\mathscr J}\) is locally complete
intersection in \(X_{\mathscr J}\).}
\begin{enumerate}
\item[\SGAInlineRefTarget{sga3:iii:proposition:4-6-5:item:a}(a)] \textit{Condition \SGARefLink{sga3:iii:proposition:4-6-5:item:a}{(a)} of \SGARefLink{sga3:iii:proposition:4-5:item:i}{4.5(i)} is then satisfied; moreover,
  every \(Y\) flat over \(S\) lifting \(Y_{\mathscr J}\) is locally
  complete intersection in \(X\).}
\item[\SGAInlineRefTarget{sga3:iii:proposition:4-6-5:item:b}(b)] \textit{If, in addition, \(Y_0\) is affine, condition \SGARefLink{sga3:iii:proposition:4-6-5:item:b}{(b)} of
  the cited result is also satisfied.}
\end{enumerate}

\noindent\textit{Proof.}
The first assertion of \SGARefLink{sga3:iii:proposition:4-6-5:item:a}{(a)} follows from \SGARefLink{sga3:iii:lemma:4-6-3}{Lemma 4.6.3}, and the second then
follows from \SGARefLink{sga3:iii:lemma:4-6-4}{Lemma 4.6.4}. Moreover, the hypothesis implies (cf. \SGARefLink{sga3:iii:definition:4-6-2}{4.6.2})
that \(\mathscr N_{Y_{\mathscr J}/X_{\mathscr J}}\) is a finite locally
free \(\mathscr O_{Y_{\mathscr J}}\)-module. Hence the
\(\mathscr O_{Y_0}\)-module
\[
 \mathscr A_0=
 \operatorname{Hom}_{\mathscr O_{Y_0}}
 \left(
  \mathscr N_{Y_{\mathscr J}/X_{\mathscr J}}
   \otimes_{\mathscr O_{Y_{\mathscr J}}}\mathscr O_{Y_0},
  \mathscr J\otimes_{\mathscr O_{S_0}}\mathscr O_{Y_0}
 \right)
\]
is quasi-coherent (cf. EGA I, 1.3.12), and therefore
\(R^1\Gamma(Y_0,\mathscr A_0)=0\) when \(Y_0\) is affine.
\hfill\(\square\)

\medskip
\noindent\SGAInlineRefTarget{sga3:iii:remark:4-6-6}\textbf{Remark 4.6.6.}\textsuperscript{\normalfont\ref{fnref:III-88}}
Let us conclude this paragraph with an example showing that, under the
hypotheses of \SGARefLink{sga3:iii:lemma:4-6-3}{Lemma 4.6.3}, a regular sequence
\((\overline g_1,\overline g_2)\) generating the ideal
\(\overline I=g(\overline E)\) need not lift to a regular sequence in
\(B\).

Let \(k\) be a field, put \(\overline A=k[X,Y]\), let
\(k\varepsilon\) denote the \(\overline A\)-module
\(\overline A/(X,Y)\), so that
\[
 P\varepsilon=P(0,0)\varepsilon
 \qquad(P\in\overline A),
\]
and let \(A=\overline A\oplus k\varepsilon\), where
\(J=k\varepsilon\) is a square-zero ideal. Thus
\(A/J=\overline A\).

The algebra \(B=A\otimes_k k[Z,T]\) is free, and hence flat, over
\(A\), and \(\overline B=k[X,Y,Z,T]\). Put
\[
 \overline g_1=XZ-YT,
 \qquad
 \overline g_2=XZ-1.
\]
Since \(\overline g_1\) is irreducible,
\(\overline B/(\overline g_1)\) is an integral domain; hence
\((\overline g_1,\overline g_2)\) is a regular sequence in
\(\overline B\), generating the ideal
\[
 \overline I=(XZ-1,YT-1).
\]
Therefore
\[
 \overline C=\overline B/\overline I
 =k[X,Y,X^{-1},Y^{-1}]
 =\overline A[X^{-1},Y^{-1}]
\]
is a flat \(\overline A\)-algebra, and also a flat \(A\)-algebra. But
every lift of \(\overline g_1\) to \(B\) has the form
\[
 XZ-YT+\lambda\varepsilon,
 \qquad \lambda\in k[Z,T],
\]
and therefore annihilates \(\varepsilon\).

\smallskip
\noindent{\footnotesize\textit{Translator's note.} The French text prints
\(XY-ZT+\lambda\varepsilon\) for this lift. Since the element being
lifted is \(\overline g_1=XZ-YT\), the displayed lift above restores
that same polynomial before adding \(\lambda\varepsilon\).}

\medskip
\noindent\SGAInlineRefTarget{sga3:iii:paragraph:4-7}\textbf{4.7.}
\SGARefLink{sga3:iii:paragraph:4-7}{Remark 4.7} has been suppressed here and placed in \SGARefLink{sga3:iii:remark:4-5-1}{4.5.1}.

